\documentclass[conference]{IEEEtran}
\usepackage{cite}
\usepackage{amsmath,amssymb,amsfonts}
\usepackage{microtype}
\usepackage{graphicx}
\usepackage{textcomp}
\usepackage{xcolor}
\usepackage{hyperref}
\usepackage{float}
\usepackage{booktabs}
\usepackage{enumitem} 
\usepackage{dblfloatfix}
\usepackage[section]{placeins}
\setlength{\textfloatsep}{8pt plus 2pt minus 2pt}
\setlength{\floatsep}{6pt plus 2pt minus 2pt}
\setlength{\intextsep}{6pt plus 2pt minus 2pt}
\begin{document}

\title{DeepfakeGuard: An Open-Source Multi-Modal Toolkit for Detecting AI-Generated Video and Audio}

\author{
\IEEEauthorblockN{Aryan Biswas, Kevin Hu, Eitan Zur, Muhammad Hamza Khan, Emmanuel Davidson}
\IEEEauthorblockA{
    \textit{Queen's University, School of Computing} \\
    Kingston, Ontario, Canada \\
    \href{mailto:21ab43@queensu.ca}{21ab43@queensu.ca},
    \href{mailto:20ch35@queensu.ca}{20ch35@queensu.ca},
    \href{mailto:22TLG3@queensu.ca}{22TLG3@queensu.ca},
    \href{mailto:23qj44@queensu.ca}{23qj44@queensu.ca},
    \href{mailto:20eobd@queensu.ca}{20eobd@queensu.ca}
}
}
\maketitle

% ================================================================
%  ABSTRACT
% ================================================================

\begin{abstract}
AI-generated video and audio have become an everyday attack surface for fraud, harassment, and political manipulation, yet effective defensive tooling remains largely confined to research labs. The 2026 International AI Safety Report documents that human observers correctly flag high-quality synthetic media only a small fraction of the time, and industry threat intelligence confirms hundreds of thousands of deepfakes already circulating on major platforms.

We present DeepfakeGuard, an open-source Python library that unifies three complementary detection modalities into a single, pip-installable framework: (1)~a DINOv3 Vision Transformer detector that achieves 0.88 AUROC under strict cross-dataset evaluation via parameter-efficient LayerNorm tuning; (2)~a LipFD audio-visual detector that identifies lip-sync deepfakes by exploiting temporal inconsistency between audio and visual lip movements; and (3)~a D3 temporal-volatility detector that requires zero training and flags AI-generated video through second-order motion features. A domain-aware ensemble system fuses all three modalities using trust-weighted scoring with applicability gating and outlier veto, and an optional post-hoc Vision-Language Model module provides human-readable forensic explanations grounded in visual evidence. All components are exposed through a unified Python API and a Streamlit demonstration interface. We argue that open, multi-modal defensive tooling is a necessary counterweight to the democratization of generative AI.
\end{abstract}

\begin{IEEEkeywords}
Deepfake detection, multi-modal forensics, Vision Transformers, lip-sync detection, training-free detection, ensemble fusion, explainable AI, AI safety
\end{IEEEkeywords}

% ================================================================
%  INTRODUCTION
% ================================================================

\section{Introduction}

Generative AI systems have crossed a critical threshold: synthetic video and audio are now sufficiently convincing that ordinary viewers cannot reliably distinguish them from authentic media~\cite{ai_safety_report}. High-quality deepfakes are already being used for executive-impersonation fraud~\cite{cnn2024arup}, non-consensual sexual imagery~\cite{sensity2024}, and political misinformation~\cite{ai_safety_report}, yet effective defensive tooling remains largely confined to research labs. Open-source generators such as DeepFaceLab~\cite{deepfacelab} and voice-cloning APIs~\cite{valle} have substantially lowered the barrier to producing manipulated media, while state-of-the-art detectors remain scattered across fragile research codebases with narrow modality coverage. Regulatory frameworks such as the EU AI Act~\cite{eu_ai_act} impose transparency obligations on synthetic media but do not prescribe detection mechanisms, leaving a tooling gap the research community must fill.

DeepfakeGuard addresses this gap by unifying three complementary detection modalities into a single, pip-installable Python library designed around four goals: (1)~\emph{multi-modal coverage} across visual, audio-visual, and temporal signals; (2)~\emph{robust generalization} via self-supervised features and training-free methods; (3)~\emph{accessible deployment} through a simple Python API; and (4)~\emph{interpretable output} via an optional VLM explanation module.

\subsection{Contributions}
This paper makes the following contributions:
\begin{enumerate}[leftmargin=*]
    \item We present \textbf{DeepfakeGuard}, an open-source Python toolkit that unifies three complementary detection modalities (visual, audio-visual, and temporal) into a single pip-installable library with a consistent Python API (\texttt{detect\_video}, \texttt{ensemble\_detect\_video}).
    \item We contribute a \textbf{parameter-efficient DINOv3 visual detector} that tunes only 8.3\% of a ViT-B/16 backbone to achieve 0.88 cross-dataset AUROC on Celeb-DF~v2, outperforming Xception by 23 AUROC points under the same cross-dataset protocol (Section~\ref{sec:DINOv3}).
    \item We \textbf{re-engineer LipFD}~\cite{lipfd2024} from a batch-oriented research training pipeline into a runtime inference library, resolving undocumented preprocessing quirks (Section~\ref{sec:lipfd}), and implement the \textbf{D3 training-free temporal detector}~\cite{zheng2025d3} end-to-end, requiring no AI-generated training data (Section~\ref{sec:d3}).
    \item We design and implement a \textbf{domain-aware ensemble fusion system} that combines all three modalities through trust-weighted scoring with applicability gating and outlier veto (Section~\ref{sec:ensemble}), and an optional \textbf{VLM semantic explainability module} that produces visual-evidence-grounded forensic explanations (Section~\ref{sec:vlm}).
    \item We analyze the \textbf{ethical implications} of deploying open deepfake detectors, including dual-use risks, demographic bias in training data, and the danger of over-reliance on automated tools (Section~\ref{sec:ethics}).
\end{enumerate}

% ================================================================
%  BACKGROUND & RELATED WORK
% ================================================================

\section{Background and Related Work}

\subsection{The Deepfake Generation Landscape}
Deepfake generation techniques have evolved rapidly across both visual and audio domains. \emph{Face-swap} methods such as DeepFaceLab~\cite{deepfacelab} transplant one person's likeness onto another's body in video, while \emph{face-reenactment} methods such as Face2Face~\cite{thies2016face2face} manipulate facial expressions while preserving identity. More recently, diffusion-based video generators such as OpenAI's Sora~\cite{sora2024} can synthesize photorealistic video from text prompts alone, making full video generation accessible to non-experts for the first time. \emph{Lip-sync} generators such as Wav2Lip~\cite{prajwal2020wav2lip} re-animate a speaker's lips to match arbitrary target audio, enabling convincing impersonation even when the face itself is unaltered. On the audio side, text-to-speech systems including VALL-E~\cite{valle} can clone a speaker's voice from seconds of reference audio, defeating traditional speaker-verification systems. The diversity of these manipulation types (face-swap, face-reenactment, lip-sync, full synthesis, and voice cloning) means that no single detection modality can cover all threat vectors, motivating the multi-modal approach taken by DeepfakeGuard.

\subsection{Visual Deepfake Detection}
Visual detection methods broadly fall into four families: artifact-based methods~\cite{face_artifacts}; temporal methods~\cite{temporal_detection}; biological signal methods~\cite{biological_signals}; and representation learning methods that train discriminative features end-to-end. Among representation-learning approaches, Xception~\cite{rossler2019faceforensics} and EfficientNet~\cite{efficientnet} have established strong baselines on FaceForensics++, but suffer AUROC degradation of 0.05--0.20 when evaluated on unseen datasets~\cite{celebdf}. The DeepfakeBench benchmark~\cite{yan2023deepfakebench}, which evaluates 15 detectors across 9 datasets under a unified protocol, confirms that this cross-dataset generalization gap remains the central challenge in visual forensics.

Recent work has explored two promising directions. First, Ojha et al.~\cite{ojha2023universal} showed that using a frozen CLIP feature space with simple nearest-neighbour classification yields surprisingly strong generalization (+15 mAP over the prior state of the art on unseen diffusion and autoregressive models), suggesting that large pretrained representations contain latent forensic signals. Second, self-supervised Vision Transformers such as DINOv2~\cite{oquab2024dinov2} attend to low-frequency structural cues rather than generator-specific artifacts~\cite{rethinking_forgery}, making them robust backbones for forensic fine-tuning. The subsequent introduction of learnable register tokens~\cite{darcet2024registers} mitigates high-norm token artifacts in DINOv2 attention maps, yielding improved feature quality; the \texttt{timm} library distributes this variant as ``DINOv3'' (\texttt{vit\_base\_patch16\_dinov3}), and we adopt this naming throughout.\footnote{We follow the naming convention used by the \texttt{timm} model library, where the DINOv2 backbone with register tokens from Darcet et al.~\cite{darcet2024registers} is distributed under the identifier \texttt{vit\_base\_patch16\_dinov3}.} Our DINOv3 detector builds on this direction, combining parameter-efficient LayerNorm tuning of a DINOv3 ViT-B/16 backbone with angular-margin metric learning to maximize cross-dataset transfer.

\subsection{Audio-Visual Lip-Sync Detection}
A complementary detection strategy exploits the tight temporal coupling between speech audio and lip movements. When a deepfake replaces a speaker's face or re-animates their lips without perfectly re-synthesizing the corresponding lip dynamics, subtle audio-visual mismatches arise that are invisible to the human eye but measurable in feature space.

Early work in this direction detected coarse audio-visual misalignment in talking-head videos~\cite{temporal_detection}, but modern lip-sync generators have closed the perceptual gap, requiring more fine-grained analysis. LipFD~\cite{lipfd2024} (\emph{Lips Are Lying}, NeurIPS 2024) formalizes this insight by combining a frozen CLIP ViT-L/14 encoder for global audio-visual features with a Region-Awareness ResNet-50 backbone that attends to multi-scale lip crops, trained on the AVLips dataset. By combining global context with fine-grained spatial attention over the lip region, LipFD captures subtle temporal inconsistencies that purely visual detectors miss. Our LipFD-based modality follows this approach and is detailed in Section~\ref{sec:lipfd}.

\subsection{Training-Free Detection}
Supervised detectors require labeled data from every new generator family, creating a perpetual update cycle. The D3 (Detection by Difference of Differences) method~\cite{zheng2025d3} (ICCV 2025) avoids this problem by observing that real video exhibits high second-order temporal volatility (reflecting Newtonian physics), while AI-generated video exhibits unusually low volatility due to the smoothness built into diffusion and autoregressive generators. Because this property is measured using any pre-trained vision encoder (xCLIP-B/16 in our implementation) without fine-tuning, D3 naturally generalizes to unseen generators and requires no AI-generated training data whatsoever. Our D3 integration is detailed in Section~\ref{sec:d3}.

% ================================================================
%  SYSTEM OVERVIEW
% ================================================================

\section{System Overview} \label{sec:system}

\subsection{Architecture}
DeepfakeGuard serves two primary audiences: researchers who need reproducible baselines, and developers who need a programmatic API for production pipelines. The library is organized around a central \texttt{DeepfakeGuard} orchestrator class that manages detector initialization, weight loading, and result aggregation. Each detection modality is encapsulated as an independent module that conforms to a shared interface: given a file path, it returns a normalized score in $[0, 1]$ (where 1 = certainly fake), a binary label, and a details dictionary containing modality-specific diagnostics.

\begin{table}[!t]
\centering
\caption{Detection modalities included in DeepfakeGuard v0.5.}
\label{tab:modalities}
\begin{tabular}{l@{\hskip 6pt}l@{\hskip 6pt}l@{\hskip 6pt}l}
\toprule
\textbf{Modality} & \textbf{Signal} & \textbf{Backbone} & \textbf{Training} \\
\midrule
DINOv3      & Face crops  & ViT-B/16         & Fine-tuned \\
LipFD       & Lips + audio & CLIP ViT-L/14 + ResNet-50 & Fine-tuned \\
D3          & Full frames & xCLIP-B/16       & Training-free \\
\bottomrule
\end{tabular}
\end{table}

At inference time, the orchestrator routes the input to the active detector, applies modality-specific preprocessing, and returns a unified result dictionary. Detectors can be switched at runtime via a single \texttt{set\_detector} call, and all three can be run jointly via the \texttt{ensemble\_detect\_video} entry-point (Section~\ref{sec:ensemble}).

\subsection{Python API}
The public interface consists of two entry-points. For single-detector use: \texttt{guard.detect\_video(path)} runs the currently active modality. For ensemble use: \texttt{ensemble\_detect\_video(guards, path)} runs all supplied detectors and produces a fused verdict. An optional VLM explainability pass is triggered by passing \texttt{vlm\_backend} and \texttt{vlm\_api\_key} to the ensemble call (Section~\ref{sec:vlm}).

% ================================================================
%  DINOv3 VISUAL DETECTOR
% ================================================================

\section{DINOv3 Visual Detector} \label{sec:DINOv3}

\subsection{Motivation}
Prior CNN-based forensic detectors can over-rely on dataset- and generator-specific cues, which often leads to degraded performance under distribution shift (e.g., unseen generators or differing compression pipelines). To improve cross-dataset robustness, we utilize self-supervised DINOv3 Vision Transformer features, which inherently attend to generalizable structural features, and fine-tune a small subset of parameters. Our approach achieves 0.88 cross-dataset AUROC while fine-tuning only 8.3\% of the model's parameters.

\subsection{Pipeline Overview}
The pipeline operates in three stages: (1) face extraction via MTCNN~\cite{mtcnn} with vertical-shift cropping; (2) frame encoding through a DINOv3 ViT-B/16 backbone with LayerNorm tuning; and (3) classification via a linear probe augmented with angular-margin metric learning. Figure~\ref{fig:architecture} illustrates the full inference pipeline.

% Full-width pipeline diagram placed at top of page
\begin{figure*}[!t]
\centering
\includegraphics[width=0.88\textwidth]{dinov3_pipeline_cropped.png}
\caption{DINOv3 visual detection pipeline. Frames are uniformly sampled from input video, faces are aligned and JPEG-simulated, then encoded by a partially fine-tuned DINOv3 backbone and classified as real or fake.}
\label{fig:architecture}
\end{figure*}

\subsection{Data Pipeline}

\noindent\textbf{Training set.} FaceForensics++ (FF++)~\cite{rossler2019faceforensics} is used exclusively for training, encompassing five manipulation methods at c23 (light) compression.

\noindent\textbf{Cross-dataset evaluation set.} Celeb-DF v2~\cite{celebdf} is used exclusively for unseen evaluation, with no fine-tuning on Celeb-DF.

\noindent\textbf{Paired sampling.} Each DataLoader batch yields 16 real and 16 fake frames from the same source video identity, producing an effective batch of 32 images with perfect class balance to reduce background shortcuts.

\noindent\textbf{Face preprocessing.} MTCNN~\cite{mtcnn} detects faces (confidence $\ge 0.95$). A square crop is extracted with 25\% padding and a vertical upward shift of 10\% of the face height to better capture lower-face artifacts.

\subsection{DINOv3 Encoder with LayerNorm Tuning}
The backbone is a \textbf{DINOv3 ViT-B/16} encoder producing 768-dimensional embeddings. Rather than fine-tuning the entire 86M-parameter model (which risks overfitting to training-set artifacts), we apply a three-stage partial tuning strategy: (1)~freeze all encoder parameters; (2)~unfreeze only the LayerNorm layers across all transformer blocks (${\sim}$40K parameters, 0.04\% of the model); (3)~additionally unfreeze the final transformer block, bringing the total to ${\sim}$7.1M trainable parameters (8.3\%). This approach allows the model to adapt its internal feature normalization and its final representation layer to the forensic task, while keeping the bulk of the pretrained visual knowledge intact.

\subsection{Classification Head and Metric Learning}
A single linear layer maps the encoder's 768-dimensional output to two class scores (real vs.\ fake). The model is trained with standard cross-entropy loss, plus an additional \emph{angular-margin} term that pushes the learned representations of real and fake faces apart. Intuitively, if two faces in the same batch have opposite labels (one real, one fake), this term penalizes any similarity between their internal representations, encouraging the model to place real and fake samples in clearly separated regions of the feature space:
\begin{equation} \label{eq:combined}
    \mathcal{L} = \mathcal{L}_{\text{CE}} + \lambda \cdot \frac{1}{|\mathcal{P}|}\sum_{(i,j)\in\mathcal{P}}\max\big(0,\ \cos(\mathbf{e}_i, \mathbf{e}_j) + m\big)
\end{equation}
where $\mathcal{P}$ is the set of all real-fake pairs in a batch, $\cos(\mathbf{e}_i, \mathbf{e}_j)$ measures how similar two embeddings are (1 = identical, $-$1 = opposite), and $m = 0.6$ is a margin that controls how far apart they must be pushed. We set the balancing weight $\lambda = 0.5$. This separation improves generalization to unseen datasets because the model learns a robust boundary between real and fake rather than memorizing dataset-specific shortcuts.

\subsection{Compression-Aware Training}
When users upload media to social platforms like Twitter or WhatsApp, the videos are heavily compressed, stripping away high-frequency data. To ensure our model does not overfit to pristine training data and fail on real-world degradation, we apply a stochastic augmentation pipeline during training:
\begin{itemize}[leftmargin=*]
    \item \textbf{JPEG compression}: $q \sim \mathcal{U}[50, 90]$, $p = 0.3$.
    \item \textbf{Gaussian blur}: $\sigma \sim \mathcal{U}[0.1, 2.0]$, $p = 0.2$.
    \item \textbf{Colour and Geometric}: Jitter, horizontal flip ($p=0.5$), rotation ($\pm 15^\circ$).
\end{itemize}

\subsection{Experiments}

\textbf{Evaluation Protocol.} Models are trained purely on FF++ and evaluated on Celeb-DF v2. We report AUROC (primary) and accuracy.

\begin{table}[!t]
\centering
\caption{Cross-dataset detection on Celeb-DF v2 (trained on FF++).}
\label{tab:results}
\begin{tabular}{lcc}
\toprule
\textbf{Method} & \textbf{Acc.} & \textbf{AUROC} \\
\midrule
MesoInception4~\cite{afchar2018mesonet}      & --   & 0.536 \\
Xception-c23~\cite{rossler2019faceforensics} & --   & 0.653 \\
Face X-ray~\cite{li2020facexray}             & --   & 0.741 \\
\textbf{DINOv3 (Ours)}  & \textbf{0.780} & \textbf{0.880} \\
\bottomrule
\end{tabular}
\end{table}

Our DINOv3 modality outperforms Xception-c23 by 23 AUROC points under the same cross-dataset protocol.

Table~\ref{tab:ablation} details the ablation study, demonstrating that our approach achieves superior generalization with far fewer tuned parameters than previous CNN baselines.

\begin{table}[!t]
\centering
\caption{Ablation on Celeb-DF v2: each row adds one component.}
\label{tab:ablation}
\begin{tabular}{lcc}
\toprule
\textbf{Configuration} & \textbf{AUROC} & \textbf{Trainable} \\
\midrule
Frozen DINOv3 + linear head         & 0.810 & ${\sim}$1.5K \\
\quad + LayerNorm tuning            & 0.845 & ${\sim}$40K  \\
\quad + Last-block unfreezing       & 0.865 & ${\sim}$7.1M \\
\quad + Metric learning ($\lambda\!=\!0.5$) & \textbf{0.880} & ${\sim}$7.1M \\
\bottomrule
\end{tabular}
\end{table}

\FloatBarrier

% ================================================================
%  AUDIO & TEMPORAL SECTIONS
% ================================================================

\section{LipFD Audio-Visual Detector} \label{sec:lipfd}

\subsection{Motivation}
The DINOv3 and D3 modalities in DeepfakeGuard operate exclusively on the visual channel.  However, a large and growing class of deepfakes, specifically lip-sync manipulations, targets the \emph{relationship} between what a person says and how their lips move.  In such forgeries the face may appear genuine, and the visual stream alone offers few useful clues.  To close this gap, we integrate the LipFD detector~\cite{lipfd2024} (\emph{Lips Are Lying}, NeurIPS 2024), which was specifically designed to exploit temporal inconsistency between audio signals and visual lip movements.  LipFD is the only modality in DeepfakeGuard that jointly reasons over audio and video, making it complementary to both DINOv3 (which analyzes spatial face artifacts) and D3 (which measures temporal motion volatility without audio).

\subsection{Architecture}
LipFD uses a dual-pathway architecture that fuses global audio-visual context with fine-grained spatial attention over the lip region.  Its design is grounded in the observation from Liu et al.~\cite{lipfd2024} that lip-sync deepfakes introduce subtle misalignments between the mel-spectrogram of the audio and the visual lip dynamics, which can be captured by combining a large-scale vision-language encoder with a region-focused convolutional backbone.

\begin{enumerate}[leftmargin=*]
    \item \textbf{Global audio-visual stream.}  A composite image is formed by vertically stacking a mel-spectrogram slice (500${\times}$2500\,px) above five consecutive, horizontally concatenated video frames (each 500${\times}$500\,px).  This 1000${\times}$2500 composite is resized to $1120{\times}1120$ and downsampled to $224{\times}224$ via a learnable $5{\times}5$ stride-5 convolution, then passed through a frozen CLIP ViT-L/14 encoder~\cite{radford2021clip} to produce a 768-dimensional global feature vector.  Because the mel-spectrogram and frames are spatially juxtaposed in a single image, the CLIP encoder implicitly learns coarse audio-visual correspondence through its pre-trained visual-linguistic attention heads.

    \item \textbf{Region-awareness stream.}  A ResNet-50 backbone~\cite{he2016resnet} independently encodes three progressively zoomed centre crops of each frame at scales 1.0${\times}$ (full 224\,px), 0.75${\times}$ (centre 168\,px), and 0.45${\times}$ (centre 102\,px, i.e.\ the lip region).  Each 2048-d regional feature is concatenated with the 768-d CLIP global vector, producing a 2816-d fused representation per scale.  A learned sigmoid gate (\texttt{nn.Linear(2816, 1)} $\to$ \texttt{Sigmoid}) assigns an importance weight to each scale, and softmax normalization ensures the weights sum to one.  The weighted combination is averaged across all five frames to yield a single fused feature vector.
\end{enumerate}

\noindent A linear classifier maps this fused vector to a binary logit.  Training uses cross-entropy loss augmented by a \emph{Region-Awareness (RA) loss}:
\begin{equation}
    \mathcal{L}_{\text{RA}} = \frac{1}{B}\sum_{b=1}^{B}\sum_{i=1}^{F} \frac{10}{\exp\bigl(\alpha_{\max}^{(b,i)} - \alpha_{0}^{(b,i)}\bigr)}
\end{equation}
where $\alpha_{\max}^{(b,i)}$ is the maximum attention weight across scales for frame $i$ of sample $b$, and $\alpha_{0}^{(b,i)}$ is the weight assigned to the full-frame (1.0${\times}$) crop.  This loss penalizes uniform attention distributions, encouraging the model to concentrate on the most informative crop, which in practice is the 0.45${\times}$ lip region.  Figure~\ref{fig:lipfd_arch} illustrates the full pipeline.

\begin{figure}[!t]
\centering
\includegraphics[width=0.88\columnwidth]{lipfd_pipeline.png}
\caption{LipFD audio-visual detection pipeline.  A composite image (mel-spectrogram + video frames) is processed by a frozen CLIP encoder and a region-aware ResNet-50, then fused via attention-weighted combination.}
\label{fig:lipfd_arch}
\end{figure}



\subsection{Integration and Implementation}\label{sec:lipfd_impl}
The original LipFD codebase~\cite{lipfd2024} is structured as a research training pipeline: an offline preprocessing script converts an entire dataset of videos into composite PNG images on disk (${\sim}$60\,GB for AVLips), which are then loaded by a PyTorch \texttt{DataLoader} during training and validation.  This design is appropriate for reproducing the paper's experiments but is not suitable for a runtime detection library that must process arbitrary video files on demand.

Our integration therefore preserves the model \emph{architecture} while re-engineering the surrounding infrastructure:

\smallskip\noindent\textbf{Model architecture (preserved).}  The neural network (the CLIP ViT-L/14 global encoder, the $5{\times}5$ stride-5 downsampling convolution, the ResNet-50 region-awareness backbone, the attention fusion mechanism, and the linear classifier) matches the original LipFD design, and we load the authors' pretrained checkpoint with compatible parameter tensor shapes (\texttt{strict=False}).  The computation graph is not modified; any difference in prediction behaviour would constitute a bug in our integration.

\smallskip\noindent\textbf{Inference preprocessing (new).}  We wrote a complete video-to-tensor pipeline that replaces the offline batch preprocessing script:
\begin{enumerate}[leftmargin=*]
    \item Ten groups of five consecutive frames are uniformly sampled from the video via OpenCV, using frame-level seeking (\texttt{CAP\_PROP\_POS\_FRAMES}) instead of decoding the entire video into memory.
    \item The audio track is extracted to a temporary WAV file via FFmpeg, then converted to a mel-spectrogram in memory using librosa.  If the video has no audio, a grey placeholder is substituted so that the visual pathway still functions.
    \item Each frame group is assembled into a composite image (mel-spectrogram above frames) in RGB, then converted to BGR via \texttt{cv2.cvtColor} to match the colour order of the original training data.
    \item Multi-scale crops (1.0${\times}$, 0.75${\times}$, 0.45${\times}$) are extracted and resized to $224{\times}224$ for the ResNet backbone.
\end{enumerate}

\smallskip\noindent\textbf{Reproduction-critical details.}  Two undocumented behaviours in the original training code required careful replication: (1)~CLIP-standard input normalization is applied but immediately overwritten by the crop extraction code, so the model was effectively trained on raw BGR pixels in $[0,\,255]$; and (2)~the crop loop uses 1-pixel horizontal offsets rather than 500-pixel frame strides.  We replicate both quirks exactly; deviating from either causes substantial accuracy degradation.  The resulting pipeline is wrapped in a \texttt{LipFDDetector} class that runs end-to-end inference in mini-batches of~4.

\subsection{Experiments}
We evaluate on FakeAVCeleb v1.2~\cite{fakeavceleb2021}, a multimodal deepfake benchmark containing real and fake videos across three manipulation types.  Each experiment uses 250 real and 250 fake videos (balanced).

\begin{table}[!t]
\centering
\caption{LipFD detection on FakeAVCeleb v1.2 (250 real + 250 fake per row).}
\label{tab:lipfd_results}
\begin{tabular}{lccccc}
\toprule
\textbf{Manipulation} & \textbf{Acc.} & \textbf{AUROC} & \textbf{Prec.} & \textbf{Rec.} & \textbf{F1} \\
\midrule
FakeVid\,+\,FakeAud   & \textbf{91.2\%} & \textbf{0.962} & 88.4\% & 94.8\% & 0.915 \\
FakeVid\,+\,RealAud   & 75.2\% & 0.821 & 83.5\% & 62.8\% & 0.717 \\
RealVid\,+\,FakeAud   & 49.6\% & 0.513 & 48.3\% & 11.6\% & 0.187 \\
\bottomrule
\end{tabular}
\end{table}

\noindent\textbf{Results.}  Table~\ref{tab:lipfd_results} reveals a clear pattern.  On its target domain (lip-synced video with synthesized audio), LipFD achieves 91.2\% accuracy and 0.962 AUROC.  Performance degrades on face-swap videos (75.2\%) and drops to chance on audio-only fakes (49.6\%), confirming that the model's discriminative power is primarily visual: the ResNet backbone detects spatial artifacts in the lip region, while the CLIP encoder provides supporting audio-visual context.  Separately, all 250 genuine videos in the real class were correctly identified as real 87.6\% of the time (TNR), yielding a false positive rate of 12.4\%.

\smallskip\noindent\textbf{Modern lip-sync evasion.}  We also tested three Wav2Lip-generated~\cite{prajwal2020wav2lip} deepfakes.  All three scored well within the ``real'' range.  One plausible explanation is dataset and generator mismatch: LipFD was trained on the AVLips dataset (which includes earlier lip-sync methods such as MakeItTalk, TalkLip, and SadTalker), while more recent systems like Wav2Lip may better preserve audio-visual consistency, evading the learned discriminative cues.

\smallskip\noindent\textbf{Complementarity.}  These results illustrate why a multi-modal ensemble is necessary.  LipFD excels at lip-sync deepfakes where DINOv3 and D3 show uncertainty (${\sim}$50\% scores), but it cannot detect face-swaps or audio-only forgeries, which are the strengths of the visual and temporal modalities respectively.  No single detector covers all threat models.

% ================================================================
% D3 TRAINING-FREE TEMPORAL DETECTOR
% ================================================================

\section{D3 Training-Free Temporal Detector} \label{sec:d3}

\subsection{Motivation}
Supervised deepfake detectors require labeled AI-generated data for every new generator, limiting their scalability and real-world applicability~\cite{zheng2025d3}. The D3 (Detection by Difference of Differences) method addresses this by exploiting a physics-grounded invariant: real-world videos exhibit complex, \emph{chaotic} second-order acceleration patterns governed by Newtonian mechanics, whereas AI-generated videos produced by diffusion and autoregressive generators display unusually smooth, constrained motion due to the smoothness built into their training process. Because this property can be measured using any pre-trained vision encoder without fine-tuning, D3 generalizes to unseen generators and requires no AI-generated training data whatsoever.

The D3 framework approximates these physical dynamics using second-order central difference approximations to measure the volatility of semantic motion:
\begin{align}
    f''(x) &= \frac{f(x+h) - 2f(x) + f(x-h)}{h^2} \\
           &= \frac{f'(x) - f'(x-h)}{h}
\end{align}
By calculating these numerical derivatives across frame embeddings, D3 identifies the ``suppressed'' motion signatures typical of current AI generators.

\subsection{Pipeline Overview}
Given a video $X \in \mathbb{R}^{T \times 3 \times H \times W}$ uniformly sampled into $T = 32$ frames at interval $\Delta t$, D3 applies a four-stage feature extraction pipeline progressing from zero-order position features to second-order acceleration and a final scalar volatility score, as illustrated in Figure~\ref{fig:d3flow}.

\begin{figure}[!t]
    \centering
    \includegraphics[width=\columnwidth]{flow2.png}
    \caption{D3 feature extraction pipeline. Real videos exhibit volatile second-order optical flow patterns (left); AI-generated videos show suppressed, constrained dynamics (right). Source:~\cite{zheng2025d3}.}
    \label{fig:d3flow}
\end{figure}

\subsection{Feature Extraction and Volatility Computation}

The D3 pipeline quantifies motion dynamics through a hierarchical extraction process. First, Zero-Order Features are captured by passing individual video frames through an xCLIP-B/16 encoder to obtain high-dimensional semantic embeddings that represent the ``position'' of objects in feature space. From these, First-Order Features are derived by calculating the rate of change between consecutive frame embeddings, effectively measuring the ``velocity'' of semantic shifts over time. 

To capture the underlying physics, the system then computes Second-Order Features, which represent the ``acceleration'' or change in velocity between frames. Finally, the Volatility score is determined by measuring the statistical variance (standard deviation) of these acceleration features. This final scalar value identifies the ``chaos'' of the motion; real videos typically yield high volatility due to natural physical complexity, while AI-generated videos yield low volatility due to artificial motion smoothing.

\subsection{Classification}
Videos are classified by comparing the volatility score against a threshold $\tau$ calibrated on real-video validation data only (no AI-generated labels required):
\begin{equation}
    \text{label} = \begin{cases} \text{Real} & \text{if } \sigma(F_2) > \tau \\ \text{Fake} & \text{if } \sigma(F_2) \leq \tau \end{cases}
\end{equation}
Threshold $\tau$ is selected to maximize F1 on a held-out set of real videos, using scikit-learn PR-curve analysis.

\subsection{Implementation}
The D3 modality is implemented entirely in Python with no training loop: videos are decoded and uniformly sampled to $T{=}32$ frames, xCLIP-B/16 embeddings (via HuggingFace Transformers) are extracted per-frame and passed through the $F_0 \to F_1 \to F_2 \to \sigma$ pipeline. The full inference path requires no GPU fine-tuning and runs on CPU for short clips.

\subsection{Experiments}

\noindent\textbf{Evaluation protocol.} We evaluate on real videos drawn from Kinetics-400 paired against AI-generated videos from Pika and Sora. We report Average Precision (AP) as the primary metric.

\begin{table}[!t]
\centering
\caption{D3 detection performance by generator.}
\label{tab:d3results}
\begin{tabular}{lcc}
\toprule
Generator & Avg.\ Volatility $\sigma$ & AP \\
\midrule
Pika  & 1.42 & 89.76\% \\
Sora  & 6.81 & $\approx$65\% \\
\bottomrule
\end{tabular}
\end{table}

\noindent\textbf{Results.} D3 achieves 89.76\% AP on Pika-generated videos. The lower AP on Sora is consistent with the observation that more realistic motion dynamics make second-order volatility cues less discriminative: Sora's higher average volatility (6.81 vs.\ 1.42 for Pika) suggests that its generated motion more closely resembles real-world physical complexity.

Figure~\ref{fig:d3heatmap} shows xCLIP attention aligned with second-order flow variance, confirming that volatility concentrates around semantically meaningful moving objects consistent with Newtonian inertia.

\begin{figure}[!b]
    \centering
    \includegraphics[width=0.5\columnwidth]{rebuttal_heatmap4.png}
    \caption{xCLIP attention heatmap aligned with second-order flow variance, confirming physics-consistent localization on moving semantics. Source:~\cite{zheng2025d3}.}
    \label{fig:d3heatmap}
\end{figure}

% ================================================================
%  ENSEMBLE DETECTION SYSTEM
% ================================================================

\section{Ensemble Detection System} \label{sec:ensemble}

Sections~\ref{sec:DINOv3}--\ref{sec:d3} establish that each detector excels on a distinct forgery family while degrading outside its domain. A na\"ive majority vote is insufficient because two inapplicable detectors can outvote an applicable one. The ensemble system addresses this through domain-aware weighted fusion, invoked via the \texttt{ensemble\_detect\_video()} entry-point.

\subsection{Applicability Gating}
Before fusion, each detector is assigned an applicability factor $a_i \in [0, 1]$ based on input-specific heuristics derived from the detector's own diagnostic output:

\begin{itemize}[leftmargin=*]
    \item \textbf{DINOv3} is scored by the number of faces detected across sampled frames: $\geq$8 frames with faces yields $a = 1.0$, 4--7 yields $a = 0.75$, and fewer yields $a = 0.5$.
    \item \textbf{D3} applicability is determined by scene motion sufficiency, computed as the ratio of the video's measured volatility to the classification threshold. Videos with very low temporal motion (ratio below 0.45) receive reduced applicability; videos above 0.75 receive $a = 1.0$.
    \item \textbf{LipFD} is gated by two criteria.  First, \emph{sample count}: $\geq$8 extractable audio-visual lip-sync samples yields $a = 1.0$, 4--7 yields $a = 0.75$, and fewer yields $a = 0.45$.  Videos that fail outright (e.g., no audio track) receive $a = 0.0$, effectively removing LipFD from the fusion.  Second, \emph{score-pattern validation}: if all per-sample scores are near-zero with negligible variance ($\sigma < 0.001$, $\max < 0.01$), LipFD receives $a = 0.15$ because such degenerate output indicates no classifiable lip-sync content was present (e.g., no visible face or no meaningful audio), and the near-zero score should not be mistaken for a confident REAL verdict.
\end{itemize}

\subsection{Outlier Veto}
If exactly one detector scores above 0.80 while all others score below 0.35, the lone detector is flagged as a likely domain mismatch. Its fusion weight is multiplied by 0.05 (a 95\% reduction). This prevents, for example, DINOv3 from dominating the verdict by detecting ``face artefacts'' in a legitimate AI-generated animation where both D3 and LipFD report low suspicion.

\subsection{Weighted Fusion}
Each detector's contribution to the final score is governed by:
\begin{equation} \label{eq:fusion}
    w_i = \pi_i \cdot (0.1 + 0.9 \cdot c_i) \cdot a_i
\end{equation}
where $\pi_i$ is a fixed trust prior reflecting each detector's overall reliability (DINOv3: 1.0, LipFD: 0.9, D3: 0.65), $c_i = 2|s_i - 0.5|$ is the certainty of detector $i$'s score (ranging from 0 at maximum ambiguity to 1 at maximum confidence), and $a_i$ is the applicability factor described above. The raw fused score is the weighted mean:
\begin{equation}
    s_{\text{raw}} = \frac{\sum_i s_i \cdot w_i}{\sum_i w_i}
\end{equation}
This raw score is then sharpened via a logistic function centred at 0.5:
\begin{equation} \label{eq:sharpen}
    s_{\text{ensemble}} = \frac{1}{1 + \exp\!\bigl(-k\,(s_{\text{raw}} - 0.5)\bigr)}, \quad k = 4
\end{equation}

The moderate sharpening constant $k = 4$ pushes scores away from the decision boundary to produce more decisive verdicts without distorting the ranking. When a VLM backend is configured, the ensemble verdict is accompanied by a VLM-generated forensic explanation (Section~\ref{sec:vlm}); otherwise a rule-based natural-language summary is provided as a fallback, covering inter-detector agreement, per-detector rationale with domain-specific caveats, and applicability warnings.

% ================================================================
%  VLM SEMANTIC EXPLAINABILITY
% ================================================================

\section{VLM Semantic Explainability} \label{sec:vlm}

To help human reviewers understand the ensemble verdict, DeepfakeGuard includes an optional post-hoc Vision-Language Model (VLM) module that provides visual-evidence-grounded explanations for both FAKE and REAL outcomes. Critically, the VLM is purely interpretive: it runs only \emph{after} all detection scores have been finalized and has zero influence on the detection decision. It degrades gracefully to \texttt{available=False} if the required backend dependencies or API keys are not present, with no impact on the detection pipeline.

\subsection{Pipeline}
Six frames are uniformly sampled from the video and stitched into a labelled $2 \times 3$ grid image. This grid, together with the ensemble score and verdict, is sent to a VLM with a structured forensic prompt. The prompt instructs the model to act as a digital forensics analyst and identify artifacts across four categories:
\begin{itemize}[leftmargin=*]
    \item \textbf{Anatomical errors:} unnatural facial geometry, impossible skin textures, asymmetric features, abnormal eye reflections.
    \item \textbf{Physics violations:} lighting inconsistent with the scene, impossible shadows, hair that defies gravity.
    \item \textbf{Temporal inconsistencies:} flickering textures across frames, identity drift, abrupt appearance changes.
    \item \textbf{AI generation artifacts:} GAN-style noise, blending seams around face boundaries, over-smoothed skin, checkerboard patterns.
\end{itemize}

The VLM response is parsed into a structured output containing a natural-language explanation, the applicable artifact categories, a self-reported confidence level (\texttt{high}/\texttt{medium}/\texttt{low}), and key frame indices identifying the strongest evidence. Robust parsing with regex fallback handles malformed JSON responses.

\subsection{Backends}
Three interchangeable backends are supported, all using identical prompts and output schemas:
\begin{itemize}[leftmargin=*]
    \item \textbf{OpenAI} (\texttt{openai}): GPT-4o mini via the OpenAI API. Cloud-based; requires an API key.
    \item \textbf{Anthropic} (\texttt{anthropic}): Claude via the Anthropic API. Cloud-based; requires an API key.
    \item \textbf{Local} (\texttt{qwen2vl}): Qwen2-VL-7B-Instruct running on-device via HuggingFace Transformers. Requires a CUDA-capable GPU; no API key needed.
\end{itemize}

\noindent Each backend is an optional dependency installed via \texttt{pip install deepfake\_guard[vlm]} (OpenAI), \texttt{[vlm-anthropic]} (Anthropic), or \texttt{[vlm-local]} (Qwen2-VL). Missing dependencies or API failures degrade gracefully with no impact on the detection pipeline.

\FloatBarrier
% ================================================================
%  ETHICS DISCUSSION
% ================================================================

\section{Ethical Considerations} \label{sec:ethics}

\subsection{Asymmetries and Justice}
The current generative-AI ecosystem exhibits a deep asymmetry: offensive tools (deepfake generators) receive vastly more investment, user adoption, and open-source community support than their defensive counterparts (deepfake detectors). Consumer face-swap applications and voice-cloning services have been downloaded or accessed by millions of users~\cite{sensity2024}, while the most widely cited detection papers typically ship as fragile research code with little support for real-world integration.

This imbalance has concrete justice implications: the people most exposed to deepfake harms (women targeted by non-consensual sexual imagery, activists in repressive regimes, low-level employees in financial institutions) are the least likely to have access to specialized forensic expertise. DeepfakeGuard is an explicit attempt to reduce this asymmetry. We caution against treating any automated detector as a complete solution or as a substitute for institutional responsibility.

\subsection{Dual-Use and Model Probing}
Any detection system can be used adversarially, for example by allowing attackers to iteratively refine deepfakes until they evade the detector. We accept this risk as inherent to security research: making detection methods openly available raises the floor for defenders, whereas obscurity primarily benefits well-resourced attackers. To mitigate dual-use risks in practice, we recommend that deployers: (1) avoid exposing raw model confidences to untrusted users; (2) rate-limit and log queries to detection APIs; and (3) combine model outputs with independent verification steps such as provenance checks and human review.

\subsection{Bias, Coverage, and Contestability}
Our DINOv3 modality is trained on FaceForensics++ and evaluated on Celeb-DF v2, datasets that over-represent Western, lighter-skinned celebrities~\cite{rossler2019faceforensics,celebdf}. The LipFD modality is trained on AVLips, which primarily contains English-language talking-head videos; its accuracy on other languages, accents, or non-frontal poses has not been systematically evaluated. Before deploying DeepfakeGuard in high-stakes environments, institutions should conduct local evaluations on demographically representative data and implement appeal or override mechanisms.

\subsection{Limitations and Future Work}
DeepfakeGuard does not address the truthfulness of verbal claims, the presence of coercion in otherwise real media, or the psychological harms inflicted on victims of synthetic imagery. Each modality also has strict input-domain requirements: DINOv3 requires a detectable face in the frame, LipFD requires both visible lips and an audio track, and D3 requires sufficient scene motion to compute meaningful volatility scores. Content that lacks these signals (e.g., a faceless screencast or a silent video) will fall outside the toolkit's coverage entirely. Beyond these scope constraints, several concrete limitations remain: (1)~LipFD cannot detect modern lip-sync generators such as Wav2Lip that closely replicate natural audio-visual synchronization (Section~\ref{sec:lipfd}); (2)~D3's volatility signal weakens as video generators improve their physics fidelity (Table~\ref{tab:d3results}); and (3)~all three modalities have been evaluated primarily on English-language, Western-demographic data.

Future research should explore tighter integration between detection and provenance tools (e.g., C2PA content credentials), systematic audits across languages and demographics, and participatory design with the communities most affected by deepfakes.

% ================================================================
%  CONCLUSION
% ================================================================

\section{Conclusion} \label{sec:conclusion}

We have presented DeepfakeGuard, an open-source multi-modal toolkit that unifies visual, audio-visual, and temporal deepfake detection into a single pip-installable Python library. By packaging a fine-tuned DINOv3 visual detector (0.88 cross-dataset AUROC), a LipFD audio-visual lip-sync detector (0.962 AUROC on lip-sync fakes), and a training-free D3 temporal detector (89.76\% AP) behind a consistent Python API, DeepfakeGuard lowers the barrier to deploying deepfake defenses from bespoke research pipelines to a three-line function call. A domain-aware ensemble system fuses all three modalities using trust-weighted scoring with applicability gating and outlier veto, while an optional VLM module provides human-readable forensic explanations grounded in visual evidence. Our experiments confirm the complementarity of these modalities: each covers threat vectors the others miss, validating the multi-modal design.

As generative AI continues to advance, we believe that open, accessible, and multi-modal detection tooling is not merely useful but ethically necessary. The toolkit, pre-trained weights, and evaluation code are publicly available at \url{https://github.com/aryanbiswas16/DeepFakeGuard}.

\section*{Author Contributions}
A.~Biswas conceived the project, designed the overall system architecture, implemented the DINOv3 visual detector, and led the writing of this paper. K.~Hu integrated the LipFD audio-visual detector (Section~\ref{sec:lipfd}). E.~Zur designed and implemented the D3 temporal detector (Section~\ref{sec:d3}). M.~H.~Khan designed and implemented the VLM semantic explainability module (Section~\ref{sec:vlm}). E.~Davidson designed and implemented the ensemble detection system (Section~\ref{sec:ensemble}).

\section*{Acknowledgment}
The authors thank Queen's University's School of Computing for supporting this research.

% ================================================================
%  REFERENCES
% ================================================================

\begin{thebibliography}{00}

\bibitem{ai_safety_report}
Y.~Bengio \emph{et al.}, ``International {AI} Safety Report 2026,'' International Scientific Report on the Safety of Advanced {AI}, Feb. 2026.
\newblock [Online]. Available: \url{https://internationalaisafetyreport.org/sites/default/files/2026-02/international-ai-safety-report-2026.pdf} (accessed Mar.~2026).

\bibitem{sensity2024}
Sensity AI, ``The State of Deepfakes 2024: Landscape, Threats, and Detection,'' Sensity AI, Amsterdam, The Netherlands, Tech. Rep., 2024.
\newblock [Online]. Available: \url{https://sensity.ai/deepfake-threat} (accessed Mar.~2026).

\bibitem{cnn2024arup}
H.~Chen and K.~Magramo, ``Finance worker pays out \$25 million after video call with deepfake `chief financial officer,'\,'' \emph{CNN}, Feb.~4, 2024.
\newblock [Online]. Available: \url{https://www.cnn.com/2024/02/04/asia/deepfake-cfo-scam-hong-kong-intl-hnk} (accessed Mar.~2026).

\bibitem{valle}
C.~Wang, S.~Chen, Y.~Wu, Z.~Zhang, L.~Zhou, S.~Liu, Z.~Chen, Y.~Liu, H.~Wang, J.~Li, L.~He, S.~Zhao, and F.~Wei, ``Neural codec language models are zero-shot text to speech synthesizers,'' \emph{arXiv preprint arXiv:2301.02111}, 2023.

\bibitem{deepfacelab}
I.~Perov, D.~Gao, N.~Chervoniy, K.~Liu, S.~Marangonda, C.~Um\'{e}, Mr.~Dpfks, C.~S.~Facez, L.~RP, J.~Jiang, S.~Zhang, P.~Wu, B.~Zhou, and W.~Zhang, ``{DeepFaceLab}: Integrated, flexible and extensible face-swapping framework,'' \emph{arXiv preprint arXiv:2005.05535}, 2020.

\bibitem{thies2016face2face}
J.~Thies, M.~Zollh\"{o}fer, M.~Stamminger, C.~Theobalt, and M.~Nie{\ss}ner, ``{Face2Face}: Real-time face capture and reenactment of {RGB} videos,'' in \emph{Proc. IEEE/CVF Conf. Comput. Vision Pattern Recognit. (CVPR)}, 2016, pp.~2387--2395.

\bibitem{face_artifacts}
Y.~Li and S.~Lyu, ``Exposing {DeepFake} videos by detecting face warping artifacts,'' in \emph{Proc. IEEE/CVF Conf. Comput. Vision Pattern Recognit. Workshops (CVPRW)}, 2019.

\bibitem{temporal_detection}
D.~G\"{u}era and E.~J.~Delp, ``Deepfake video detection using recurrent neural networks,'' in \emph{Proc. IEEE Int. Conf. Adv. Video Signal-Based Surveillance (AVSS)}, 2018, pp.~1--6.

\bibitem{biological_signals}
U.~A.~Ciftci, I.~Demir, and L.~Yin, ``{FakeCatcher}: Detection of synthetic portrait videos using biological signals,'' \emph{IEEE Trans. Pattern Anal. Mach. Intell.}, vol.~42, no.~11, pp.~2726--2740, 2020.

\bibitem{rossler2019faceforensics}
A.~R\"{o}ssler, D.~Cozzolino, L.~Verdoliva, C.~Riess, J.~Thies, and M.~Nie{\ss}ner, ``{FaceForensics++}: Learning to detect manipulated facial images,'' in \emph{Proc. IEEE/CVF Int. Conf. Comput. Vision (ICCV)}, 2019, pp.~1--11.

\bibitem{efficientnet}
M.~Tan and Q.~V.~Le, ``{EfficientNet}: Rethinking model scaling for convolutional neural networks,'' in \emph{Proc. Int. Conf. Mach. Learn. (ICML)}, 2019, pp.~6105--6114.

\bibitem{celebdf}
Y.~Li, X.~Yang, P.~Sun, H.~Qi, and S.~Lyu, ``{Celeb-DF}: A large-scale challenging dataset for {DeepFake} forensics,'' in \emph{Proc. IEEE/CVF Conf. Comput. Vision Pattern Recognit. (CVPR)}, 2020, pp.~3207--3216.

\bibitem{oquab2024dinov2}
M.~Oquab, T.~Darcet, T.~Moutakanni, H.~V.~Vo, M.~Szafraniec, V.~Khalidov, P.~Fernandez, D.~Haziza, F.~Massa, A.~El-Nouby, M.~Assran, N.~Ballas, W.~Galuba, R.~Howes, P.-Y.~Huang, S.-W.~Li, I.~Misra, M.~Rabbat, V.~Sharma, G.~Synnaeve, H.~Xu, H.~J\'{e}gou, J.~Mairal, P.~Labatut, A.~Joulin, and P.~Bojanowski, ``{DINOv2}: Learning robust visual features without supervision,'' \emph{Trans. Mach. Learn. Res. (TMLR)}, 2024.

\bibitem{darcet2024registers}
T.~Darcet, M.~Oquab, J.~Mairal, and P.~Bojanowski, ``Vision {Transformers} need registers,'' in \emph{Proc. Int. Conf. Learn. Representations (ICLR)}, 2024.

\bibitem{rethinking_forgery}
R.~Shao, T.~Wu, and Z.~Liu, ``Detecting and grounding multi-modal media manipulation,'' in \emph{Proc. IEEE/CVF Conf. Comput. Vision Pattern Recognit. (CVPR)}, 2023.

\bibitem{lipfd2024}
W.~Liu, J.~Lin, B.~Li, S.~Peng, J.~Fang, and Z.~He, ``Lips are lying: Spotting the temporal inconsistency between audio and visual in lip-syncing deepfakes,'' in \emph{Advances in Neural Inf. Process. Syst. (NeurIPS)}, 2024.

\bibitem{zheng2025d3}
E.~Zheng, T.~Kang, and H.~Zou, ``{D3}: Training-free {AI}-generated video detection using second-order features,'' in \emph{Proc. IEEE/CVF Int. Conf. Comput. Vision (ICCV)}, 2025.

\bibitem{mtcnn}
K.~Zhang, Z.~Zhang, Z.~Li, and Y.~Qiao, ``Joint face detection and alignment using multitask cascaded convolutional networks,'' \emph{IEEE Signal Process. Lett.}, vol.~23, no.~10, pp.~1499--1503, 2016.

\bibitem{li2020facexray}
L.~Li, J.~Bao, T.~Zhang, H.~Yang, D.~Chen, F.~Wen, and B.~Guo, ``Face {X-Ray} for more general face forgery detection,'' in \emph{Proc. IEEE/CVF Conf. Comput. Vision Pattern Recognit. (CVPR)}, 2020, pp.~5001--5010.

\bibitem{afchar2018mesonet}
D.~Afchar, V.~Nozick, J.~Yamagishi, and I.~Echizen, ``{MesoNet}: A compact facial video forgery detection network,'' in \emph{Proc. IEEE Int. Workshop Inf. Forensics Security (WIFS)}, 2018, pp.~1--7.

\bibitem{fakeavceleb2021}
H.~Khalid, S.~Kim, S.~Tariq, and S.~S.~Woo, ``{FakeAVCeleb}: A novel audio-video multimodal deepfake dataset,'' in \emph{Advances in Neural Inf. Process. Syst. (NeurIPS) Datasets Benchmarks Track}, 2021.

\bibitem{prajwal2020wav2lip}
K.~R.~Prajwal, R.~Mukhopadhyay, V.~P.~Namboodiri, and C.~V.~Jawahar, ``A lip sync expert is all you need for speech to lip generation in the wild,'' in \emph{Proc. ACM Int. Conf. Multimedia (MM)}, 2020, pp.~484--492.

\bibitem{radford2021clip}
A.~Radford, J.~W.~Kim, C.~Hallacy, A.~Ramesh, G.~Goh, S.~Agarwal, G.~Sastry, A.~Askell, P.~Mishkin, J.~Clark, G.~Krueger, and I.~Sutskever, ``Learning transferable visual models from natural language supervision,'' in \emph{Proc. Int. Conf. Mach. Learn. (ICML)}, vol.~139, 2021, pp.~8748--8763.

\bibitem{he2016resnet}
K.~He, X.~Zhang, S.~Ren, and J.~Sun, ``Deep residual learning for image recognition,'' in \emph{Proc. IEEE/CVF Conf. Comput. Vision Pattern Recognit. (CVPR)}, 2016, pp.~770--778.

\bibitem{eu_ai_act}
European Parliament and Council of the European Union, ``Regulation ({EU}) 2024/1689 laying down harmonised rules on artificial intelligence ({AI Act}),'' \emph{Official Journal of the European Union}, vol.~L, 2024/1689, Jul. 2024.
\newblock [Online]. Available: \url{http://data.europa.eu/eli/reg/2024/1689/oj} (accessed Mar.~2026).

\bibitem{ojha2023universal}
U.~Ojha, Y.~Li, and Y.~J.~Lee, ``Towards universal fake image detectors that generalize across generative models,'' in \emph{Proc. IEEE/CVF Conf. Comput. Vision Pattern Recognit. (CVPR)}, 2023.

\bibitem{sora2024}
OpenAI, ``Video generation models as world simulators,'' Tech. Rep., Feb. 2024.
\newblock [Online]. Available: \url{https://openai.com/index/video-generation-models-as-world-simulators} (accessed Mar.~2026).

\bibitem{yan2023deepfakebench}
Z.~Yan, Y.~Zhang, X.~Yuan, S.~Lyu, and B.~Wu, ``{DeepfakeBench}: A comprehensive benchmark of deepfake detection,'' in \emph{Advances in Neural Inf. Process. Syst. (NeurIPS)}, 2023.

\end{thebibliography}

\end{document}
